%% toc: P15 160 160.9
\chapter{Second Derivatives, Sine Products and the Weak-Law Derivative Criterion}
\label{sec:p15-w02c-p15-013-second-derivatives-and-products}

Three statements in Durrett's Section~3.3 are taken up here in the form in
which the source states them. The first is the sentence that introduces his
Theorem~3.3.21: existence of a second derivative of a characteristic
function at the origin forces a finite second moment
\cite[p.~136, Theorem~3.3.21]{DurrettPTEAuthor2019}. The theorem itself
assumes only that the upper limit of a symmetric second difference quotient
is finite; that form is \cref{thm:p15-characteristic-second-moment}. This
chapter proves the derivative reading, including the step from a second
derivative to convergence of the quotient. The second is the derivation that
Exercise~3.3.10 sketches before it asks for a characteristic-function proof:
repeated use of the double-angle identity gives the infinite cosine product
for \((\sin t)/t\) \cite[p.~133, Exercise~3.3.10]{DurrettPTEAuthor2019}. The
third is the attribution with which Durrett closes Exercise~3.3.14
\cite[p.~136, Exercise~3.3.14]{DurrettPTEAuthor2019}. The mathematics of that
exercise is \cref{thm:p15-criteria-pitman-law}.

\paragraph{Inputs and conclusions.}
The first section takes a finite real random variable \(X\) with
characteristic function \(\varphi\) (\cref{def:p15-characteristic-function})
and assumes that \(\varphi''(0)\) exists. It concludes that
\(\mathbb EX^2<\infty\) and \(\varphi''(0)=-\mathbb EX^2\), and that \(\varphi\)
is then twice continuously differentiable on all of \(\mathbb R\). For a
characteristic function, existence of \(\varphi''(0)\) turns out to be
equivalent to each of the following: Durrett's upper-limit hypothesis,
convergence of the symmetric quotient, and a finite second moment. The
second section takes a real number \(t\) and proves the identity
\((\sin t)/t=\prod_{m\geq1}\cos(t/2^m)\) using only trigonometry. It shows
that the partial products do not converge uniformly on \(\mathbb R\). It
then uses the identity to prove that the binary series of fair signs is
uniform on \((-1,1)\), which runs the source's characteristic-function
interpretation in the reverse direction.

\section{A second derivative at the origin gives a second moment}
\label{sec:p15-w02c-p15-013-second-derivative}

Throughout this section \(X\) is a finite real random variable with Borel
probability law \(\mu\) and characteristic function
\(\varphi(t)=\int e^{itx}\,\mu(dx)\), as in
\cref{sec:p15-characteristic-moments}.

\begin{hbconvention}{def:p15-w02c-p15-013-second-derivative}{Second derivative at the origin, the symmetric quotient and the real part}
Derivatives of the complex function \(\varphi\) of a real variable are taken
componentwise (\cref{def:p15-characteristic-complex-conventions}).
\emph{\(\varphi''(0)\) exists} means: there is \(\delta>0\) such that
\(\varphi\) is differentiable at every point of \((-\delta,\delta)\), and the
derivative \(\varphi'\) is differentiable at \(0\); then
\[
 \varphi''(0)=\lim_{h\to0}\frac{\varphi'(h)-\varphi'(0)}{h}\in\mathbb C .
\]
Nothing is assumed about \(\varphi''\) away from \(0\).
The symmetric quotient of \cref{thm:p15-characteristic-second-moment} is kept:
\[
 Q(h)=\frac{\varphi(h)-2\varphi(0)+\varphi(-h)}{h^2}\qquad(h\neq0).
\]
Write \(u(t)=\operatorname{Re}\varphi(t)\) for the real part of \(\varphi\).
Three hypotheses on \(\varphi\) at the origin are compared:
\begin{enumerate}
\item[(D)] \(\varphi''(0)\) exists;
\item[(S)] the limit \(\lim_{h\to0}Q(h)\) exists and is finite;
\item[(U)] \(\limsup_{h\downarrow0}Q(h)>-\infty\), which is the hypothesis
      of Durrett's Theorem~3.3.21 \cite[p.~136]{DurrettPTEAuthor2019}.
\end{enumerate}
\end{hbconvention}

\begin{hblemma}{thm:p15-w02c-p15-013-quotient-limit}{proved=annex:p15-w02c-p15-013-quotient-limit}{A second derivative is the limit of the symmetric second difference}
Let \(f\) be a real function that is differentiable at every point of an
open interval \((-\delta,\delta)\), \(\delta>0\), and whose derivative \(f'\)
is differentiable at \(0\). Then
\[
 \lim_{h\to0}\frac{f(h)-2f(0)+f(-h)}{h^2}=f''(0).
\]
\end{hblemma}

\begin{hbtheorem}{thm:p15-w02c-p15-013-second-derivative-moment}{proved=annex:p15-w02c-p15-013-second-derivative-moment}{Durrett's reading of Theorem 3.3.21: existence of a second derivative}
Suppose \(\varphi''(0)\) exists in the sense of
\cref{def:p15-w02c-p15-013-second-derivative}. Then:
\begin{enumerate}
\item \(Q(h)\to\varphi''(0)\) as \(h\to0\); in particular
      \(\limsup_{h\downarrow0}Q(h)=\varphi''(0)>-\infty\), so the
      hypothesis of \cref{thm:p15-characteristic-second-moment} holds;
\item \(\mathbb EX^2<\infty\), and
      \(\varphi''(0)=-\mathbb EX^2\); in particular \(\varphi''(0)\) is real
      and at most \(0\);
\item \(\mathbb E|X|<\infty\), \(\varphi'(0)=i\,\mathbb EX\), and \(\varphi\)
      has continuous derivatives through order \(2\) on all of
      \(\mathbb R\), with
      \(\varphi''(t)=-\mathbb E[X^2e^{itX}]\) for every real \(t\).
\end{enumerate}
\end{hbtheorem}

\begin{hbcorollary}{thm:p15-w02c-p15-013-four-conditions}{proved=annex:p15-w02c-p15-013-four-conditions}{Four equivalent second-order conditions for a characteristic function}
For the characteristic function \(\varphi\) of a finite real random variable
\(X\), the conditions (D), (S), (U) of
\cref{def:p15-w02c-p15-013-second-derivative} and the condition
\(\mathbb EX^2<\infty\) are equivalent. When they hold,
\[
 \varphi''(0)=\lim_{h\to0}Q(h)=-\mathbb EX^2 .
\]
The equivalence of (U) with (D) uses that \(\varphi\) is a characteristic
function. For the real function \(g(t)=|t|\), the upper limit in (U) is
\(+\infty\), yet \(g\) is not differentiable at \(0\).
\end{hbcorollary}

\begin{hbexample}{def:p15-w02c-p15-013-uniform-second-moment}{The uniform law on \((-1,1)\)}
Let \(X\) be uniform on \((-1,1)\). By
\cref{thm:p15-characteristic-elementary-examples} with \(c=1\),
\(\varphi(t)=(\sin t)/t\), with \(\varphi(0)=1\). Dividing the sine series
by \(t\) gives, for every real \(t\),
\[
 \varphi(t)=\sum_{k=0}^{\infty}\frac{(-1)^k t^{2k}}{(2k+1)!}
           =1-\frac{t^2}{6}+\frac{t^4}{120}-\cdots,
\]
a power series that converges on all of \(\mathbb R\). It may be
differentiated term by term, so \(\varphi''(0)\) exists and equals
\(2\cdot(-1/6)=-1/3\). By
\cref{thm:p15-w02c-p15-013-second-derivative-moment},
\(\mathbb EX^2=1/3\) and \(Q(h)=2((\sin h)/h-1)/h^2\to-1/3\). The direct
computation \(\int_{-1}^{1}x^2\cdot\tfrac12\,dx=\tfrac13\) gives the same
value.
\end{hbexample}

\begin{hbtheorem}{thm:p15-w02c-p15-013-order-contrast}{proved=annex:p15-w02c-p15-013-order-contrast}{A first derivative without a second: the symmetric stable transforms}
Let \(1<\alpha<2\) and let \(X\) have characteristic function
\(\varphi_\alpha(t)=e^{-|t|^\alpha}\), which is a characteristic function by
\cref{thm:p15-criteria-stable-functions}. Then \(\varphi_\alpha\) is
continuously differentiable on \(\mathbb R\) with \(\varphi_\alpha'(0)=0\).
For \(0<|h|\leq1\), its symmetric quotient satisfies
\(Q(h)\leq-|h|^{\alpha-2}\). Hence \(Q(h)\to-\infty\), \(\varphi_\alpha''(0)\)
does not exist, and \(\mathbb EX^2=\infty\).
\end{hbtheorem}

\begin{hbwarning}{def:p15-w02c-p15-013-order-one}{The first-order analogue fails}
At order one, existence of the derivative does not give the moment.
\Cref{thm:p15-criteria-pitman-law} exhibits a law for which
\(\varphi'(0)\) exists although \(\mathbb E|X|=\infty\). The order-two
argument rests on a sign: by conjugation symmetry,
\(Q(h)=-\int 2(1-\cos hx)h^{-2}\,\mu(dx)\), an integral of a nonpositive
integrand, so Fatou's lemma applies. The first-order quotient
\((\varphi(h)-\varphi(-h))/(2h)=i\int(\sin hx)/h\,\mu(dx)\) has an integrand
of both signs, and Fatou's lemma gives no bound from it. In the other
direction, \cref{thm:p15-w02c-p15-013-order-contrast} shows that a
continuous first derivative at the origin does not give a second derivative
there.
\end{hbwarning}

\section{The sine product from the double-angle identity}
\label{sec:p15-w02c-p15-013-sine-product}

The source's Exercise~3.3.10 says that repeated use of
\(\sin t=2\sin(t/2)\cos(t/2)\) leads to
\((\sin t)/t=\prod_{m=1}^\infty\cos(t/2^m)\), and then asks for a proof by
reading both sides as characteristic functions
\cite[p.~133, Exercise~3.3.10]{DurrettPTEAuthor2019}. The
characteristic-function proof is
\cref{thm:p15-criteria-series-examples}. The trigonometric derivation is
written out here. For real \(t\) and integers \(n\geq1\) put
\[
 P_n(t)=\prod_{m=1}^{n}\cos(t/2^m).
\]
The infinite product means \(\lim_{n\to\infty}P_n(t)\), as in
\cref{thm:p15-criteria-limit-operations}; zero limits are allowed.

\begin{hbtheorem}{thm:p15-w02c-p15-013-sine-product}{proved=annex:p15-w02c-p15-013-sine-product}{The cosine product by repeated halving}
For every real \(t\) and every integer \(n\geq1\),
\[
 \sin t=2^n\sin(t/2^n)\,P_n(t).
\]
If \(t\neq0\) and \(2^n>|t|/\pi\), then \(\sin(t/2^n)\neq0\) and
\(P_n(t)=\sin t/(2^n\sin(t/2^n))\). Consequently
\[
 \frac{\sin t}{t}=\prod_{m=1}^{\infty}\cos(t/2^m)\qquad(t\in\mathbb R),
\]
where the left side is \(1\) at \(t=0\). The convergence is not uniform on
\(\mathbb R\): for every \(n\geq1\), \(P_n(2^n\pi)=-1\), while the limit
function vanishes at \(2^n\pi\), so
\(\sup_{t\in\mathbb R}|P_n(t)-(\sin t)/t|\geq1\).
\end{hbtheorem}

\begin{hbcorollary}{thm:p15-w02c-p15-013-half-pi}{proved=annex:p15-w02c-p15-013-half-pi}{The value at \(\pi/2\) as nested square roots}
Put \(c_1=\sqrt2/2\) and \(c_{m+1}=\sqrt{(1+c_m)/2}\) for \(m\geq1\). Then
\(c_m=\cos(\pi/2^{m+1})\) for every \(m\geq1\), and
\[
 \frac{2}{\pi}=\prod_{m=1}^{\infty}c_m
 =\frac{\sqrt2}{2}\cdot\frac{\sqrt{2+\sqrt2}}{2}\cdot
  \frac{\sqrt{2+\sqrt{2+\sqrt2}}}{2}\cdots .
\]
\end{hbcorollary}

\begin{hbcorollary}{thm:p15-w02c-p15-013-binary-uniform}{proved=annex:p15-w02c-p15-013-binary-uniform}{The identity read backwards: the binary series is uniform}
Let \(\varepsilon_1,\varepsilon_2,\dots\) be independent with
\(\mathbb P(\varepsilon_j=1)=\mathbb P(\varepsilon_j=-1)=1/2\). The series
\(V=\sum_{j\geq1}\varepsilon_j2^{-j}\) converges at every outcome, its partial
sums have characteristic functions \(P_n\), and \(V\) is uniform on
\((-1,1)\). The proof uses \cref{thm:p15-w02c-p15-013-sine-product} together
with the infinite-product rule and uniqueness of laws. It does not use the
grid-counting argument of \cref{thm:p15-criteria-series-examples}.
\end{hbcorollary}

The partial products \(P_n\) are the characteristic functions of the partial
sums of \(V\). These sums converge almost surely, so their laws converge
weakly, and \cref{thm:p15-criteria-tight-transforms} gives uniform
convergence of \(P_n\) on every compact interval. The last assertion of
\cref{thm:p15-w02c-p15-013-sine-product} shows that the convergence fails
to be uniform on the whole line. This is a second instance, besides
\(\delta_{1/n}\Rightarrow\delta_0\), of the failure asked for in Exercise~3.3.13(iii)
\cite[p.~136]{DurrettPTEAuthor2019}.

\section{Attribution of the derivative criterion for the weak law}
\label{sec:p15-w02c-p15-013-pitman}

Durrett's Exercise~3.3.14 has three parts
\cite[p.~136, Exercise~3.3.14]{DurrettPTEAuthor2019}. For i.i.d.\ \(X_j\)
with characteristic function \(\varphi\) and \(S_n=X_1+\cdots+X_n\):
part~(i) derives \(S_n/n\to a\) in probability from \(\varphi'(0)=ia\);
part~(ii) derives \(\varphi(t/n)^n\to e^{iat}\) along the integers from the
convergence in probability; and part~(iii) uses part~(ii) and uniform
continuity of \(\varphi\) to reach the derivative quotient
\((\varphi(h)-1)/h\) as \(h\) decreases to \(0\) through all positive reals.
Together they give the criterion that the weak law holds exactly when
\(\varphi'(0)\) exists.

Durrett credits this criterion to E.~J.~G. Pitman (1956). He adds that John
Walsh pointed out that part~(iii) should be proved
\cite[p.~136]{DurrettPTEAuthor2019}. Durrett's reference list gives the
paper as Pitman, ``On derivatives of characteristic functions at the
origin'', \emph{Ann. Math. Statist.} 27, 1156--1160
\cite[p.~481]{DurrettPTEAuthor2019}. The step Walsh identified is needed
because part~(ii) controls \(\varphi\) only at the points \(t/n\) for a
fixed \(t\), whereas a derivative is a limit over every small \(h\). In
\cref{annex:p15.criteria-pitman-law} this step sets \(n=\lfloor1/h\rfloor\)
and \(t_n=nh\in(0,1]\), and it uses the uniform convergence on \([0,1]\)
together with a continuous logarithm. Negative \(h\) are then handled by
conjugation symmetry.

The source prints the limit in part~(iii) as \(-ia\). The value \(ia\) is
the one consistent with part~(i) and with \(\varphi'(0)=ia\), since
\(\varphi(h)=1+iah+o(h)\) is what \(\varphi'(0)=ia\) means; this is the
value proved in \cref{thm:p15-criteria-pitman-law}. The presentation there
follows Durrett's exercise route. Pitman's own paper is not reproduced, and
no statement here is drawn from it.

\section{Relationships and sources}
\label{sec:p15-w02c-p15-013-relationships}

\Cref{thm:p15-w02c-p15-013-second-derivative-moment} reduces to
\cref{thm:p15-characteristic-second-moment} (Durrett's Theorem~3.3.21,
p.~136) through \cref{thm:p15-w02c-p15-013-quotient-limit}. It then
uses \cref{thm:p15-characteristic-moment-derivatives} (Theorem~3.3.18,
p.~134) for the derivative formulas. The converse direction of
\cref{thm:p15-w02c-p15-013-four-conditions} is Theorem~3.3.18 with \(n=2\).
\Cref{thm:p15-w02c-p15-013-order-contrast} draws on the stable
transforms of \cref{thm:p15-criteria-stable-functions}.
\Cref{def:p15-w02c-p15-013-order-one} draws on the example in
\cref{thm:p15-criteria-pitman-law}.
The sine product uses the remainder bound of
\cref{thm:p15-characteristic-taylor-remainder} (Lemma~3.3.19, p.~134) for
\(\sin x/x\to1\). Its probabilistic reading is
\cref{thm:p15-criteria-series-examples}, and
\cref{thm:p15-w02c-p15-013-binary-uniform} runs that reading in reverse with
\cref{thm:p15-criteria-limit-operations} and
\cref{thm:p15-characteristic-inversion-uniqueness}.
The source locators are Durrett, \emph{Probability: Theory and Examples},
fifth edition: Exercise~3.3.10 on p.~133; Theorem~3.3.18 and
Lemma~3.3.19 on p.~134; Theorem~3.3.21, the sentence introducing it, and
Exercise~3.3.14 with its attribution on p.~136; and the Pitman entry of the
reference list on p.~481 \cite{DurrettPTEAuthor2019}.
